% !TeX encoding = UTF-8
\documentclass[variant=apuntes,cover=institutional,mode=screen]{ua-engineering-v6-article}
% !TeX encoding = UTF-8
\UASetUniversity{Universidad Austral}
\UASetFaculty{Facultad de Ingeniería}
\UASetDepartment{Departamento de Computación}
\UASetCampus{Pilar}
\UASetCareer{Ingeniería Informática}
\UASetCommission{Comisión A}
\UASetProfessor{Equipo docente}
\UASetSemester{Segundo cuatrimestre 2026}
\UASetCourse{Sistemas Distribuidos}
\UASetDate{2026}


\UASetClassNumber{05}
\UASetClassTitle{Transacciones distribuidas}
\UASetDocKind{Apuntes}

\title{Clase 05 --- Transacciones distribuidas y serializabilidad}
\author{\UAProfessor}

\date{\UADate}
\begin{document}
\maketitle

\section{Introducción}

Las transacciones permiten agrupar operaciones de forma que el sistema preserve invariantes incluso bajo concurrencia.

En sistemas distribuidos, esto implica:
\begin{itemize}
\item múltiples nodos
\item múltiples copias
\item latencias y fallas
\end{itemize}

\section{Modelo de transacciones}

Una transacción $T$ es una secuencia:
\[
T = \{r(x), w(x), commit\}
\]

Conflictos:
\begin{itemize}
\item Read-Write (RW)
\item Write-Read (WR)
\item Write-Write (WW)
\end{itemize}

\section{Serializabilidad}

\subsection{Definición}
Una ejecución es serializable si existe una ejecución secuencial equivalente.

\subsection{Grafo de precedencia}

Nodos: transacciones  
Aristas: dependencias por conflicto

\paragraph{Teorema}
El schedule es serializable ⇔ el grafo es acíclico.

\section{Anomalías}

\subsection{Dirty read}
Lectura de datos no confirmados.

\subsection{Non-repeatable read}
Lecturas inconsistentes dentro de la misma transacción.

\subsection{Phantom}
Cambios en conjuntos de resultados.

\subsection{Write skew}
Dos transacciones leen el mismo estado y escriben sin conflicto directo.

\paragraph{Ejemplo}
Dos médicos verifican guardia:
\begin{itemize}
\item ambos leen que hay otro médico
\item ambos se retiran
\end{itemize}
Resultado: ningún médico.

\section{Snapshot Isolation}

Cada transacción observa un snapshot consistente.

\subsection{Propiedad}
Evita:
\begin{itemize}
\item dirty read
\item non-repeatable read
\end{itemize}

No evita:
\begin{itemize}
\item write skew
\end{itemize}

\section{Commit distribuido}

\subsection{2PC}

Fases:
\begin{enumerate}
\item Prepare
\item Commit
\end{enumerate}

\subsection{Problema}
Bloqueo si el coordinador falla.

\section{3PC}

Introduce fase:
\begin{itemize}
\item pre-commit
\end{itemize}

Reduce bloqueo pero requiere sincronía.

\section{Relación con consenso}

\begin{itemize}
\item consenso $\Rightarrow$ orden global
\item transacciones $\Rightarrow$ consistencia lógica
\end{itemize}

\section{Conclusión}

\begin{quote}
La serializabilidad garantiza correctitud, pero requiere coordinación fuerte y tiene costo.
\end{quote}


\section{Síntesis razonada de la clase}
Esta unidad debe leerse como parte de una cadena de decisiones de diseño y no como un catálogo de mecanismos. Para estudiar el tema con profundidad conviene reconstruir cuatro preguntas: \emph{qué problema aparece}, \emph{qué garantía se desea}, \emph{qué mecanismo la aproxima} y \emph{qué costo introduce}. En cada ejemplo debe identificarse además el modelo de falla implícito.

\begin{uakeybox}
Una respuesta técnicamente madura no termina al nombrar un algoritmo o patrón: declara sus supuestos, la propiedad que preserva y el escenario en el que deja de ser suficiente.
\end{uakeybox}

\section{Bibliografía guiada y ruta de profundización}
Para profundizar esta clase se recomienda: Kleppmann, cap. 7; Gray y Reuter; Berenson et al. sobre isolation levels.

La lectura debe hacerse con preguntas concretas: (1) ¿qué modelo de sistema asume el autor?, (2) ¿qué propiedad demuestra o persigue?, (3) ¿qué falla o anomalía motiva el mecanismo?, y (4) ¿qué trade-off queda abierto?

\section{Conexión con el Trabajo Final Integrador}
El grupo debe señalar qué parte de su sistema utiliza los conceptos de esta unidad, justificar por qué son necesarios y documentar una alternativa descartada. Cuando el contenido todavía no corresponda a la etapa vigente, debe registrarse como hipótesis o decisión pendiente.

\section{Autoevaluación conceptual}
\begin{enumerate}
\item ¿Cuál es el problema distribuido concreto que motiva esta clase?
\item ¿Qué propiedad de seguridad y qué propiedad de progreso están en juego?
\item ¿Qué supuesto de red, tiempo o falla resulta decisivo?
\item ¿Qué trade-off aceptaría en un sistema real y cuál no aceptaría en un dominio crítico?
\item ¿Cómo observaría experimentalmente que la solución se comporta como espera?
\end{enumerate}

\section{Profundización autocontenida v9}
\subsection{Problema, límite e invariante}
El núcleo de esta clase es ACID, serializabilidad, grafos de precedencia, anomalías, Snapshot Isolation, 2PC y sagas. Para estudiarlo de manera autónoma conviene separar tres planos. Primero, el \emph{problema observable}: qué comportamiento incorrecto puede ver un cliente o un operador. Segundo, el \emph{límite del modelo}: qué información, sincronía o coordinación no está disponible. Tercero, el \emph{invariante}: qué propiedad no debe violarse aunque el sistema pierda progreso temporalmente.

El modelo de fallas relevante incluye concurrencia, aborto parcial, coordinador caído y violación de invariantes. Una solución debe describirse como protocolo y no solo como nombre de tecnología: actores, estados, mensajes, condición de éxito, condición de aborto y estado visible después de una interrupción. La evidencia esperada es schedule, grafo de precedencia, estado prepared y resultado observable.

\subsection{Método de análisis}
Ante un caso nuevo, siga esta secuencia:
\begin{enumerate}
\item Identifique actores, dato o recurso compartido y frontera de confianza.
\item Escriba una ejecución nominal y una ejecución adversarial.
\item Distinga seguridad (lo malo no ocurre) de progreso (algo bueno finalmente ocurre).
\item Declare qué supuesto permite avanzar: mayoría, deadline, sincronía parcial, operación idempotente, almacenamiento durable u otro.
\item Compare una alternativa más simple y una más fuerte; registre qué métrica o experimento haría cambiar la decisión.
\end{enumerate}

\subsection{Caso transferible}
Considere una plataforma de reservas que recibe operaciones duplicadas, pierde conectividad entre regiones y debe sostener un camino crítico sin ocultar el estado real al usuario. Aplique el mecanismo de esta clase a una única decisión concreta. Una respuesta madura no intenta resolver todo: delimita la operación, formula la garantía y explica la degradación esperada cuando el supuesto central deja de cumplirse.

\subsection{Errores de razonamiento frecuentes}
\begin{itemize}
\item confundir una propiedad con el producto que suele implementarla;
\item describir el camino feliz sin recorrer una falla intermedia;
\item afirmar ``exactly-once'', ``alta disponibilidad'' o ``consistencia'' sin definir alcance observable;
\item medir solo throughput aunque la hipótesis trate sobre semántica o recuperación;
\item presentar la complejidad añadida como beneficio sin compararla con la alternativa más simple.
\end{itemize}

\section{Lectura guiada}
Kleppmann, capítulo 7; Berenson et al. sobre isolation levels; Gray y Reuter.

Durante la lectura, registre: modelo de sistema, propiedad perseguida, contraejemplo que motiva el mecanismo, costo principal y una condición bajo la cual no lo elegiría. Estas cinco anotaciones convierten la bibliografía en una herramienta de diseño y no en una lista para memorizar.

\section{Aplicación al Trabajo Final Integrador}
El equipo debe producir un ADR breve con la decisión asociada a esta clase. Debe contener: contexto, dos alternativas, supuesto decisivo, garantía elegida, trade-off, evidencia pendiente y fecha de revisión. La evidencia mínima esperada para esta unidad es: schedule, grafo de precedencia, estado prepared y resultado observable.

\section{Autoevaluación avanzada}
\begin{enumerate}
\item ¿Puedo explicar la clase sin nombrar primero una tecnología?
\item ¿Puedo construir una ejecución que viole la garantía si el mecanismo se configura mal?
\item ¿Puedo distinguir seguridad, progreso y experiencia del usuario?
\item ¿Puedo proponer una medición que valide la propiedad, no solo el rendimiento?
\item ¿Puedo defender cuándo una solución más simple sería suficiente?
\end{enumerate}

\end{document}
